% GENERATED FILE. Edit canonical lessons, then run npm run generate:books.
% canonical-source-hash: fnv1a64:6b4a9f3a423d3795
% canonical-lessons: ZH-C323-shaonian, ZH-C323-shiqi, ZH-C323-niandai, ZH-W323-dai-era, ZH-C323-gudai, ZH-C323-xiandai, ZH-R323-review-318-323

\chapter{Decades and Eras}
\label{ch:zh-decades-and-eras}
\hlchaptermodality{323}

\begin{chapteropening}
\textbf{By the end of this chapter:} \emph{I can talk about a teenager, a period, a decade, ancient times and modern times.}

Complete the last lesson of chapter 323: I can talk about a teenager, a period, a decade, ancient times and modern times.
\end{chapteropening}

\section[\texorpdfstring{\zh{少年} (shàonián)}{少年 (shàonián)}]{\zh{少年} (shàonián) — teenage boy; teenager}

\label{lesson:ZH-C323-shaonian}



\begin{quote}
Before the new one: say the Mandarin for all one's life, then the Mandarin for to look back fondly on.
\end{quote}

\subsection*{You'll want to know: \zh{少年}}
\textbf{\zh{少年}} — \emph{shàonián} — \textquotedblleft{}teenage boy; teenager\textquotedblright{}.

\begin{quote}
\textbf{\zh{他十五岁的时候是一个快乐的少年}} \emph{tā shíwǔ suì de shíhou shì yí ge kuàilè de shàonián} — At fifteen he was a happy teenager.
\end{quote}

\begin{cousinweb}
\textbf{\zh{少年}} (\emph{shàonián}) is built from \zh{少} 'young' and \zh{年} 'a year', both inherited from Middle Chinese, \emph{syewH} and \emph{nen} (Guangyun, in Baxter's transcription), which glosses the first, in this falling tone, 'young'; in its rising tone, \emph{syewX}, it is 'few'.
\end{cousinweb}

\begin{grammarlens}[title={What you've built}]
One more word for telling what happened.
\end{grammarlens}

\subsection*{Your turn}
\begin{itemize}
\raggedright
  \item \emph{Say it:} \emph{shàonián}
  \item \emph{Say it:} \emph{shàonián}, once more
  \item \emph{Say it:} \emph{tā shíwǔ suì de shíhou shì yí ge kuàilè de shàonián} — At fifteen he was a happy teenager.
  \item \emph{Recall:} say the Mandarin for all one's life, then the Mandarin for to look back fondly on, then say \emph{shàonián} again
\end{itemize}

\begin{tcolorbox}[breakable,colback=teal!4,colframe=teal!35!black,title={Before you move on}]
What does \zh{少年} mean? (\textquotedblleft{}Teenage boy; teenager\textquotedblright{}.) Say it once more.
\end{tcolorbox}
\section[\texorpdfstring{\zh{时期} (shíqī)}{时期 (shíqī)}]{\zh{时期} (shíqī) — period; phase}

\label{lesson:ZH-C323-shiqi}



\begin{quote}
Before the new one: say the Mandarin for to look back fondly on, then the Mandarin for a teenage boy.
\end{quote}

\subsection*{You'll want to know: \zh{时期}}
\textbf{\zh{时期}} — \emph{shíqī} — \textquotedblleft{}period; phase\textquotedblright{}.

\begin{quote}
\textbf{\zh{那个时期当地的人都没有手机}} \emph{nàge shíqī dāngdì de rén dōu méiyǒu shǒujī} — In that period nobody there had a mobile phone.
\end{quote}

\begin{cousinweb}
\textbf{\zh{时期}} (\emph{shíqī}) is built from \zh{时} 'a time' and \zh{期} 'a set time', both inherited from Middle Chinese, \emph{dzyi} and \emph{gi} (Guangyun, in Baxter's transcription), which glosses the second 'to meet' and 'a limit'. Middle Chinese \emph{gi} would regularly give a rising tone, qí; the level tone of qī is irregular.
\end{cousinweb}

\begin{grammarlens}[title={What you've built}]
One more word for telling what happened.
\end{grammarlens}

\subsection*{Your turn}
\begin{itemize}
\raggedright
  \item \emph{Say it:} \emph{shíqī}
  \item \emph{Say it:} \emph{shíqī}, once more
  \item \emph{Say it:} \emph{nàge shíqī dāngdì de rén dōu méiyǒu shǒujī} — In that period nobody there had a mobile phone.
  \item \emph{Recall:} say the Mandarin for to look back fondly on, then the Mandarin for a teenage boy, then say \emph{shíqī} again
\end{itemize}

\begin{tcolorbox}[breakable,colback=teal!4,colframe=teal!35!black,title={Before you move on}]
What does \zh{时期} mean? (\textquotedblleft{}Period; phase\textquotedblright{}.) Say it once more.
\end{tcolorbox}
\section[\texorpdfstring{\zh{年代} (niándài)}{年代 (niándài)}]{\zh{年代} (niándài) — decade; era}

\label{lesson:ZH-C323-niandai}



\begin{quote}
Before the new one: say the Mandarin for a teenage boy, then the Mandarin for a period in history.
\end{quote}

\subsection*{You'll want to know: \zh{年代}}
\textbf{\zh{年代}} — \emph{niándài} — \textquotedblleft{}decade; era\textquotedblright{}.

\begin{quote}
\textbf{\zh{他八十年代去过中国}} \emph{tā bāshí niándài qù guo Zhōngguó} — He went to China in the eighties.
\end{quote}

The new character, in pieces: \zh{代} — five strokes: the person piece \zh{亻} on the left; on the right a bar rising to the right, a long stroke that slants down to the right and flicks up, and a dot. The next lesson puts it in your hand.

\begin{cousinweb}
\textbf{\zh{年代}} (\emph{niándài}) is built from \zh{年} 'a year' and \emph{dài} 'an age', both inherited from Middle Chinese, \emph{nen} and \emph{dojH} (Guangyun, in Baxter's transcription), which glosses the second 'to take turns' and gives this very word as its example.
\end{cousinweb}

\begin{grammarlens}[title={What you've built}]
One more word for telling what happened.
\end{grammarlens}

\subsection*{Your turn}
\begin{itemize}
\raggedright
  \item \emph{Say it:} \emph{niándài}
  \item \emph{Say it:} \emph{niándài}, once more
  \item \emph{Say it:} \emph{tā bāshí niándài qù guo Zhōngguó} — He went to China in the eighties.
  \item \emph{Recall:} say the Mandarin for a teenage boy, then the Mandarin for a period in history, then say \emph{niándài} again
\end{itemize}

\begin{tcolorbox}[breakable,colback=teal!4,colframe=teal!35!black,title={Before you move on}]
What does \zh{年代} mean? (\textquotedblleft{}Decade; era\textquotedblright{}.) Say it once more.
\end{tcolorbox}
\section[\texorpdfstring{\zh{代} (dài)}{代 (dài)}]{\zh{代} — the person piece, then a slanting hook}

\label{lesson:ZH-W323-dai-era}



\begin{quote}
The word first, then the character inside it.

\begin{itemize}
\raggedright
  \item \emph{Recall:} say \emph{niándài} — \textbf{R1}, one lesson back, before it decays
\end{itemize}

\begin{quote}
\zh{信} \zh{你}
\end{quote}

\emph{Read it:} each one aloud, with nothing said first — \emph{xìn}, \emph{nǐ}
\end{quote}

\subsection*{Script — the person piece, then a slanting hook}
\hlblockfigure{\detokenize{figures/ZH-W323-dai-era-filmstrip.pdf}}{How \zh{代} is written, stroke by stroke}

\begin{quote}
\zh{代}
\end{quote}

\textbf{Two pieces, left before right.}

\textbf{The left piece, \zh{亻}, two strokes,} as in \zh{你} and \zh{信}.

\textbf{The right piece, three strokes.} A short bar that rises a little to the right. Then a long stroke from the top, crossing the bar and slanting down to the right, flicking up at the foot. Last, a dot at the top right.

2 + 3 = five.

\subsection*{Your turn}
\begin{itemize}
\raggedright
  \item \emph{Write it:} \zh{代}, all five, with the model in front of you
  \item \emph{Cover:} write it again from memory, and count the pieces, not the strokes
\end{itemize}

\begin{tcolorbox}[breakable,colback=teal!4,colframe=teal!35!black,title={Before you move on}]
How many strokes in \zh{代}? (\textbf{Five}.) Which stroke flicks up at the foot? (\textbf{The long slanting stroke, the fourth}.) What comes last? (\textbf{A dot at the top right}.)

Source: \href{https://raw.githubusercontent.com/chanind/hanzi-writer-data/68d10a4b21150cae5e1ebbd223eed289cf32d90c/data/\%E4\%BB\%A3.json}{Hanzi Writer Data \zh{代}, pinned stroke paths and medians}.
\end{tcolorbox}
\section[\texorpdfstring{\zh{古代} (gǔdài)}{古代 (gǔdài)}]{\zh{古代} (gǔdài) — ancient times}

\label{lesson:ZH-C323-gudai}



\begin{quote}
Before the new one: say the Mandarin for a period in history, then the Mandarin for a decade of a century.
\end{quote}

\subsection*{You'll want to know: \zh{古代}}
\textbf{\zh{古代}} — \emph{gǔdài} — \textquotedblleft{}ancient times\textquotedblright{}.

\begin{quote}
\textbf{\zh{我喜欢看古代的小说}} \emph{wǒ xǐhuan kàn gǔdài de xiǎoshuō} — I like reading novels from ancient times.
\end{quote}

\begin{cousinweb}
\textbf{\zh{古代}} (\emph{gǔdài}) is built from \zh{古} 'ancient' and \zh{代} 'an age', both inherited from Middle Chinese, \emph{kuX} and \emph{dojH} (Guangyun, in Baxter's transcription), which glosses the first 'former' and the second 'to take turns'.
\end{cousinweb}

\begin{grammarlens}[title={What you've built}]
One more word for telling what happened.
\end{grammarlens}

\subsection*{Your turn}
\begin{itemize}
\raggedright
  \item \emph{Say it:} \emph{gǔdài}
  \item \emph{Say it:} \emph{gǔdài}, once more
  \item \emph{Say it:} \emph{wǒ xǐhuan kàn gǔdài de xiǎoshuō} — I like reading novels from ancient times.
  \item \emph{Recall:} say the Mandarin for a period in history, then the Mandarin for a decade of a century, then say \emph{gǔdài} again
\end{itemize}

\begin{tcolorbox}[breakable,colback=teal!4,colframe=teal!35!black,title={Before you move on}]
What does \zh{古代} mean? (\textquotedblleft{}Ancient times\textquotedblright{}.) Say it once more.
\end{tcolorbox}
\section[\texorpdfstring{\zh{现代} (xiàndài)}{现代 (xiàndài)}]{\zh{现代} (xiàndài) — modern times; modern}

\label{lesson:ZH-C323-xiandai}



\begin{quote}
Before the new one: say the Mandarin for a decade of a century, then the Mandarin for ancient times.
\end{quote}

\subsection*{You'll want to know: \zh{现代}}
\textbf{\zh{现代}} — \emph{xiàndài} — \textquotedblleft{}modern times; modern\textquotedblright{}.

\begin{quote}
\textbf{\zh{现代的人都有手机}} \emph{xiàndài de rén dōu yǒu shǒujī} — Everyone has a mobile phone these days.
\end{quote}

\begin{cousinweb}
\textbf{\zh{现代}} (\emph{xiàndài}) is built from \zh{现} 'to appear, now' and \zh{代} 'an age', both inherited from Middle Chinese, \emph{henH} and \emph{dojH} (Guangyun, in Baxter's transcription), which glosses the first 'to show'; the sense 'now' is not given there.
\end{cousinweb}

\begin{grammarlens}[title={What you've built}]
One more word for telling what happened.
\end{grammarlens}

\subsection*{Your turn}
\begin{itemize}
\raggedright
  \item \emph{Say it:} \emph{xiàndài}
  \item \emph{Say it:} \emph{xiàndài}, once more
  \item \emph{Say it:} \emph{xiàndài de rén dōu yǒu shǒujī} — Everyone has a mobile phone these days.
  \item \emph{Recall:} say the Mandarin for a decade of a century, then the Mandarin for ancient times, then say \emph{xiàndài} again
\end{itemize}

\begin{tcolorbox}[breakable,colback=teal!4,colframe=teal!35!black,title={Before you move on}]
What does \zh{现代} mean? (\textquotedblleft{}Modern times; modern\textquotedblright{}.) Say it once more.
\end{tcolorbox}
\section[(dialogue)]{Looking back: from Letters from Far Away to Decades and Eras}

\label{lesson:ZH-R323-review-318-323}



\begin{quote}
Say the Mandarin for a pen pal and an envelope.
\end{quote}

\subsection*{Your turn}
\begin{itemize}
\raggedright
  \item \emph{Say it:} a pen pal, an envelope, a postcard, a reply to a letter, to pass on a message
  \item \emph{Say it:} to return a call, confidence that it will work, an incoming call, to write to each other, phone charges
  \item \emph{Say it:} to end a call, to dial a number, to get through on the phone, to bring up a subject, a voice message
  \item \emph{Say it:} the history of a country, the people of long ago, an old book from long ago, age-old, a script for a language
  \item \emph{Say it:} in those years long ago, in one's early years, in one's later years, all one's life, to look back fondly on
  \item \emph{Say it:} a teenage boy, a period in history, a decade of a century, ancient times, modern times
  \item \emph{Read it:} \zh{信} \zh{古} \zh{代} — each one aloud, then a word that holds it
\end{itemize}

\begin{tcolorbox}[breakable,colback=teal!4,colframe=teal!35!black,title={Before you move on}]
A pen pal? (\emph{bǐyǒu}.) An envelope? (\emph{xìnfēng}.) A postcard? (\emph{míngxìnpiàn}.) A reply to a letter? (\textbf{\zh{回音}}, \emph{huíyīn}.) To pass on a message? (\emph{zhuǎngào}.) To return a call? (\textbf{\zh{回电}}, \emph{huídiàn}.) Confidence that it will work? (\textbf{\zh{信心}}, \emph{xìnxīn}.) An incoming call? (\textbf{\zh{来电}}, \emph{láidiàn}.) To write to each other? (\emph{tōngxìn}.) Phone charges? (\emph{huàfèi}.) To end a call? (\emph{guàduàn}.) To dial a number? (\emph{bōdǎ}.) To get through on the phone? (\emph{dǎtōng}.) To bring up a subject? (\textbf{\zh{说起}}, \emph{shuōqǐ}.) A voice message? (\textbf{\zh{语音}}, \emph{yǔyīn}.) The history of a country? (\emph{lìshǐ}.) The people of long ago? (\textbf{\zh{古人}}, \emph{gǔrén}.) An old book from long ago? (\textbf{\zh{古书}}, \emph{gǔshū}.) Age-old? (\textbf{\zh{古老}}, \emph{gǔlǎo}.) A script for a language? (\textbf{\zh{文字}}, \emph{wénzì}.) In those years long ago? (\textbf{\zh{当年}}, \emph{dāngnián}.) In one's early years? (\textbf{\zh{早年}}, \emph{zǎonián}.) In one's later years? (\textbf{\zh{晚年}}, \emph{wǎnnián}.) All one's life? (\textbf{\zh{一生}}, \emph{yìshēng}.) To look back fondly on? (\emph{huáiniàn}.) A teenage boy? (\textbf{\zh{少年}}, \emph{shàonián}.) A period in history? (\textbf{\zh{时期}}, \emph{shíqī}.) A decade of a century? (\textbf{\zh{年代}}, \emph{niándài}.) Ancient times? (\textbf{\zh{古代}}, \emph{gǔdài}.) Modern times? (\textbf{\zh{现代}}, \emph{xiàndài}.)
\end{tcolorbox}
